%% ma: L12-B06-0073
%% lop: 12
%% chuong: 2
%% bai: 6
%% dang: 12.6.13
%% loai: TN4
%% muc: TH
%% dap_an: "C"
%% nguon: "CĐ2.1.pdf giáo viên gửi 27/10/2026 (M07, Lê Minh Khôi – Nguyễn Anh Khôi, Chương II Chủ đề 1 Vectơ trong không gian), Câu 33"
%% kien_thuc: "Bài 6 (bình phương vô hướng, góc giữa hai vectơ)"
%% trang_thai: chinh-thuc
\begin{ex}
Cho hai vectơ $\vec a$ và $\vec b$ có $|\vec a|=5$, $|\vec b|=12$ và $\left|\vec a+\vec b\right|=13$. Đặt $\vec u=\vec a+\vec b$ và $\vec v=\vec a-\vec b$. Khi đó giá trị cosin của góc $\left(\vec u,\vec v\right)$ bằng
\choice{$\cos\left(\vec u,\vec v\right)=\dfrac{12}{13}$}{$\cos\left(\vec u,\vec v\right)=\dfrac5{12}$}{\True $\cos\left(\vec u,\vec v\right)=-\dfrac{119}{169}$}{$\cos\left(\vec u,\vec v\right)=\dfrac{119}{169}$}
\loigiai{$\left|\vec a+\vec b\right|^2=25+144+2\vec a\cdot\vec b=169\Rightarrow\vec a\cdot\vec b=0$, nên $|\vec v|=13$.

$\vec u\cdot\vec v=|\vec a|^2-|\vec b|^2=-119\Rightarrow\cos\left(\vec u,\vec v\right)=\dfrac{-119}{13\cdot13}=-\dfrac{119}{169}$.}
\end{ex}
